% ===================================================================
%  ТАБЛИЦА № 12 — Вокзал. — Железные дороги. — Суда.
%  Traduction 2026 d'apres texto/io, controlee sur texto/fr, texto/en
%  et texto/es. Consigne : voir texto/ru/10-tablica-01.tex.
% ===================================================================

%%K t12-apar-1 apar
\VUsaut{4.0mm}
\VUtitre{15.0pt}{60}{ТАБЛИЦА № 12}
\VUsaut{2.4mm}
\VUpk{11.4pt}{40}{Вокзал. --- Железные дороги. --- Суда.}
\VUsaut{3.4mm}
\textsc{Примечание.} --- \textit{Расписания в каждой стране свои, поэтому
за образец мы взяли часы} \VUgras{французской таблицы}.

%%K t12-01-1 p
1. --- Я только что совершил поездку по Германии. Перед отъездом я купил
\VUgras{расписание} \textsuperscript{(15)}, чтобы узнать часы отхода поездов.
Убедившись, что самый удобный поезд уходит из Парижа в девять вечера и
приходит в Страсбург в час тринадцать, я собрался в дорогу, а назавтра
велел отвезти себя на вокзал.

%%K t12-02-1 p
2. --- Когда я вошёл в большой зал отправления вслед за старым сельским
\VUgras{священником} \textsuperscript{(2)}, несшим свою \VUgras{дорожную сумку}
\textsuperscript{(3)}, немало \VUgras{пассажиров} \textsuperscript{(1)} уже собралось.

%%K t12-02-2 p
Так как мне надо было послать \VUgras{телеграмму} \textsuperscript{(5)}, я попросил
\VUgras{контролёра} \textsuperscript{(6)} показать мне \VUgras{телеграф}
\textsuperscript{(4)}.

%%K t12-03-1 p
3. --- Прежде чем пройти в \VUgras{зал ожидания} \textsuperscript{(7)}, я пошёл
взять обратный \VUgras{билет} \textsuperscript{(9)}.

%%K t12-04-1 p
4. --- У \VUgras{кассы} \textsuperscript{(8)} я ждал недолго, потому что народу
было мало. \VUgras{Солдат} \textsuperscript{(10)}, ехавший в отпуск, любезно
уступил мне свою очередь, так что у меня хватило времени задать несколько
вопросов \VUgras{начальнику вокзала} \textsuperscript{(13)}, который беседовал с
\VUgras{полицейским приставом} \textsuperscript{(12)} и с дежурным
\VUgras{жандармом} \textsuperscript{(11)}.

%%K t12-tit-1 sub
\VUpk{10.2pt}{40}{Сдача багажа.}
\VUsaut{3.0mm}

%%K t12-05-1 p
5. --- Купив газету в \VUgras{киоске} \textsuperscript{(14)}, я велел взвесить свой
\VUgras{багаж} \textsuperscript{(17)} --- \VUgras{носильщик} \textsuperscript{(16)} только
что привёз его на \VUgras{багажной тележке} \textsuperscript{(19)}.
\VUgras{Служащий} \textsuperscript{(22)} при \VUgras{весах} \textsuperscript{(21)} сказал
мне, что мой сундук весит тридцать пять килограммов, то есть на пять
килограммов \VUgras{сверх нормы}, за которые пришлось доплатить в
\VUgras{багажной кассе} \textsuperscript{(23)}.

%%K t12-06-1 p
6. --- Так как я приехал заранее, у меня было время посмотреть на движение
пассажиров по вокзалу. Одни сдавали или получали небольшие сумки в
\VUgras{камере хранения} \textsuperscript{(25)}; другие справлялись по
\VUgras{расписаниям} \textsuperscript{(26)}. Приезжие спешили к \VUgras{выходу}
\textsuperscript{(32)} с \VUgras{чемоданами} \textsuperscript{(24)} в руках. Некоторые
ждали у \VUgras{выдачи багажа} \textsuperscript{(27)} и подавали свою
\VUgras{квитанцию} \textsuperscript{(29)}, чтобы получить сундуки, \VUgras{ящики}
\textsuperscript{(28)} или \VUgras{корзины} \textsuperscript{(30)}.

%%K t12-07-1 p
7. --- \VUgras{Акцизные чиновники} \textsuperscript{(33)} внимательно осматривали
свёртки, нет ли чего к объявлению.

%%K t12-08-1 p
8. --- Иные пассажиры направлялись в \VUgras{буфет} \textsuperscript{(34)}, где
\VUgras{официант} \textsuperscript{(35)} усердно их обслуживал. Приходилось
спешить, потому что \VUgras{поезд} \textsuperscript{(36)}, курьерский, стоял на
вокзале недолго.

%%K t12-09-1 p
9. --- Затем я вышел на перрон отправления и стал искать прямой
\VUgras{вагон} \textsuperscript{(37)}, чтобы не пересаживаться в пути. Я вошёл в
вагон с коридором, где было \VUgras{купе} \textsuperscript{(38)} для курящих и
одно, отведённое «для дам, едущих одних».

%%K t12-tit-2 sub
\VUpk{10.2pt}{40}{Ночной отъезд.}
\VUsaut{3.0mm}

%%K t12-10-1 p
10. --- Поднявшись на \VUgras{подножку} \textsuperscript{(40)}, затворив
\VUgras{дверцу} \textsuperscript{(39)}, подняв \VUgras{штору} \textsuperscript{(41)} и
положив свои сумки в \VUgras{сетку} \textsuperscript{(43)}, я сел на
\VUgras{диван} \textsuperscript{(42)}. Увидев, как \VUgras{фонарщик} \textsuperscript{(44)}
зажигает лампы, я вспомнил, что мне ночевать в вагоне, и позвал служителя,
который отдаёт напрокат \VUgras{подушки} \textsuperscript{(45)}, чтобы взять одну.

%%K t12-11-1 p
11. --- Кондуктор пришёл проверять билеты. Я спросил его, отчего мы не
трогаемся, ведь время подошло. Он ответил, что \VUgras{обер-кондуктор}
\textsuperscript{(46)} занят погрузкой велосипедов в \VUgras{багажный вагон}
\textsuperscript{(47)}. К тому же не все \VUgras{грелки} \textsuperscript{(48)} ещё
переменили.

%%K t12-12-1 p
12. --- Но отправляться нам предстояло скоро, потому что \VUgras{кочегар}
\textsuperscript{(50)} на своём \VUgras{тендере} \textsuperscript{(49)} набирал уголь,
чтобы поддержать огонь \VUgras{паровоза} \textsuperscript{(54)}, а
\VUgras{машинист} \textsuperscript{(51)} ждал только знака от \VUgras{помощника
начальника вокзала} \textsuperscript{(52)}, стоявшего на \VUgras{перроне}
\textsuperscript{(53)}.

%%K t12-13-1 p
13. --- \VUgras{Котёл} \textsuperscript{(56)} паровоза глухо гудел; \VUgras{труба}
\textsuperscript{(55)} выбрасывала потоки дыма пополам с \VUgras{паром}
\textsuperscript{(60)}. Раздался пронзительный \VUgras{свисток} \textsuperscript{(59)}, и
\VUgras{поршни} \textsuperscript{(58)} пришли в движение, приводя в ход
\VUgras{колёса} \textsuperscript{(57)} тяжёлой машины.

%%K t12-tit-3 sub
\VUsaut{3.0mm}
\VUpk{10.2pt}{40}{Поехали!}
\VUsaut{3.0mm}

%%K t12-14-1 p
14. --- Поезд, бегущий по \VUgras{рельсам} \textsuperscript{(64)}, набирал ход;
\VUgras{перроны} \textsuperscript{(65)} кончились; мы миновали \VUgras{уборные}
\textsuperscript{(66)}, а также ряд вагонов, стоявших на другом \VUgras{пути}
\textsuperscript{(63)}: \VUgras{спальные вагоны} \textsuperscript{(67)},
\VUgras{вагон-ресторан} \textsuperscript{(68)}, \VUgras{почтовый вагон}
\textsuperscript{(69)}.

%%K t12-noto-1 noto
\VUnotes{143.2mm}{%
\textit{(*) Примечание.} --- \textit{Поездка, разумеется, описана так,
как она проходила по довоенной границе.}%
}

%%K t12-15-1 p
15. --- Затем перед нашими глазами прошла \VUgras{товарная станция}
\textsuperscript{(71)}. Под навесами служащие грузили \VUgras{товары}
\textsuperscript{(72)}, отправляемые малой скоростью. \VUgras{Сцепщики}
\textsuperscript{(76)} возле \VUgras{водонапорной башни} \textsuperscript{(73)}
передвигали \VUgras{вагоны для скота} \textsuperscript{(75)} и \VUgras{платформы}
\textsuperscript{(79)}, составляя \VUgras{товарный поезд} \textsuperscript{(80)}.

%%K t12-16-1 p
16. --- Мы миновали последнюю \VUgras{стрелку} \textsuperscript{(78)}, возле
которой стоял стрелочник; мы прошли \VUgras{семафор} \textsuperscript{(86)}, и
вокзал скрылся вдали.

%%K t12-17-1 p
17. --- Вскоре мы выехали на \VUgras{железный мост} \textsuperscript{(74)}, который
пересекает \VUgras{реку} \textsuperscript{(81)}. Миновав мост, мы побежали вдоль
реки, по которой мощные \VUgras{буксиры} \textsuperscript{(82)} тянули тяжёлые
\VUgras{баржи} \textsuperscript{(83)}. Мы видели и \VUgras{пассажирский пароход}
\textsuperscript{(84)}, только что подходивший к \VUgras{пристани} \textsuperscript{(85)}.

%%K t12-18-1 p
18. --- Промчавшись через несколько станций, мы прошли несколько
\VUgras{туннелей} \textsuperscript{(87)} и оставили позади бесчисленные \VUgras{домики}
\textsuperscript{(88)} \VUgras{путевых сторожей}. Убаюканный качкой поезда, я
заснул крепким сном.

%%K t12-tit-4 sub
\VUsaut{3.0mm}
\VUpk{10.2pt}{40}{Станция Дойч-Аврикур.}
\VUsaut{3.0mm}

%%K t12-19-1 p
19. --- Меня вдруг разбудил голос служащего, кричавшего:
«Дойч-Аврикур! \textsuperscript{(*)} Пересадка!» Это был таможенный досмотр и все
докучные пограничные формальности. Мне пришлось поставить часы по
вокзальным \VUgras{часам} \textsuperscript{(89)}, которые шли на пятьдесят пять
минут вперёд; спросонок я едва отличал \VUgras{большую стрелку}
\textsuperscript{(90)} от \VUgras{маленькой} \textsuperscript{(91)}; сам
\VUgras{циферблат} \textsuperscript{(92)} был совсем размыт утренним туманом.

%%K t12-20-1 p
20. --- Пока таможенники вели свой досмотр, я, зевая, разглядывал
\VUgras{афиши} \textsuperscript{(93)}, которыми украшены стены вокзала. Я
позавтракал, чтобы скоротать время, справился по карте германской
\VUgras{сети} \textsuperscript{(94)}, далеко ли ещё до Страсбурга, и вернулся в
свой вагон, ободрённый надеждой скоро добраться до цели.
\VUsaut{6.0mm}
\VUfilet{22mm}
